\cotefile{157-3}
\scanpage{327}
% slanted
\margin{Cf pages 11~\unsure{ff|ss} pour \unsure{les} contreexemples-clef}

\uline{\unsure{Prop\textsuperscript{n}|Prop}} (?). Soit $X$ dans $\mathrm{Cat}$, \unsure{’un|un|ou} $b,c$ dans $X$ tels que

a) $b\times c$ est dans $X$

b) $\forall x$ dans $X$, $b(x)\ (\overset{\text{déf}}{=}\operatorname{Hom}_{X}(x,b))$ et $c(x)$ sont $\ne\emptyset$, i.e. $X_{/b}\longrightarrow X$ et $X_{/c}\longrightarrow X$ sont surj. sur les objets.

Alors $X$ est $W_\infty$-asphérique.

\margin{Faux, cf plus bas. $X$ est 1-connexe, mais pas néc. 2-connexe.}

Le fait que $b$ soit un plus grand élément de $\mathrm{Ob}\,\unsure{\alpha|X}$ assure déjà la 0-connexité.

Pour la 1-connexité, prouvons que dans le groupoïde fondamental $\Pi(X)$, $a\overset{\text{déf}}{=}b\times c$ devient un élément final. Pour ceci, prouvons \unsure{d'abord|s'abord} que $\forall x$ dans $X$, \deleted{\unsure{dans}} les éléments de $a(x)=b(x)\times c(x)$ sont égaux dans $\mathrm{Fl}\,\Pi(X)$. En effet, comme $b\times c\longrightarrow b$ est iso \inserted{ou \unsure{inverso|verso}} dans $\Pi(X)$, \unsure{cel|cela}

\diagram{\begin{tikzcd}[ampersand replacement=\&]
 b\times c \arrow[r]\arrow[d] \& b\\
 c \&
\end{tikzcd}}

montre déjà que \deleted{\unsure{soit}} si $\beta,\beta'\in b(x)$, $\gamma,\gamma'\in c(x)$, \deleted{$(\unsure{\beta},\unsure{\beta}),(\unsure{\beta'},\unsure{\gamma'})$} \inserted{\unsure{d'où} $(\beta',\gamma')$} $\in b(x)\times c(x)$, on a $(\beta,\gamma)=(\beta,\gamma')$ dans $\Pi(X)$ \deleted{\unsure{aussi}} (puisque $\mathrm{pr}_1(\beta,\gamma)=\mathrm{pr}_1(\beta,\gamma')$, et de $\unsure{=|\text{même}}$ (en utilisant que $\mathrm{pr}_2$ est iso dans \unsure{même}) $(\beta,\gamma')=(\beta',\gamma')$.

Donc tous les $(\beta,\gamma)$ \deleted{\unsure{deviennent}} \inserted{$\in b(x)\times c(x)=a(x)$} sont égaux dans $\mathrm{Fl}\,\Pi(X)$. Pour achever de prouver que $a$ est objet final de $\Pi(X)$, il suffit de voir que toute flèche \uline{dans} $\Pi(X)$ d'un $x\in\mathrm{ob}\,X$ dans $a$, provient d'une flèche $x\longrightarrow a$ dans $X$.

Mais une \deleted{telle} flèche \inserted{de $\Pi(X)$} est donnée par un chemin \deleted{dans} $X$

\diagram{\begin{tikzcd}[ampersand replacement=\&,column sep=small]
 x=x_0\arrow[r,"u_1"]\arrow[dr,"\alpha_0"'] \& x_1\arrow[d,"\alpha_1"] \& x_2\arrow[l,"u_2"']\arrow[r]\arrow[dl,"\alpha_2"] \& \dots \& x_{2N}=a\arrow[l,"u_{2N}"]\\
 \& a \& \& \&
\end{tikzcd}}

\deleted{\unsure{en}} prenons une flèche $\alpha_0:x\longrightarrow a_0$ dans $X$, je dis que \unsure{la|les}

% archive sheet number: 327
